% B-59, 15.09.2026: datierte Ledger-Nachauswertung, Stützung und Reichweite.
% Beleg: thesis-evidence/w2/ledger-repetitions-20260915/ (Originale unverändert).
% M19 (12.09.2026): Leserführung und präzise Ergebnisdarstellung;
% Messwerte und historische Nachweisgrenzen bleiben erhalten.
% ============================================================
% Tabelle 7.3-1: Kernzahlen beider Fälle (Gold v2) — v02 (08.09.)
% In Overleaf ablegen als: tables/Chap7-v02/chap7.3table1.tex
% Aufruf im Fließtext: % B-59, 15.09.2026: datierte Ledger-Nachauswertung, Stützung und Reichweite.
% Beleg: thesis-evidence/w2/ledger-repetitions-20260915/ (Originale unverändert).
% M19 (12.09.2026): Leserführung und präzise Ergebnisdarstellung;
% Messwerte und historische Nachweisgrenzen bleiben erhalten.
% ============================================================
% Tabelle 7.3-1: Kernzahlen beider Fälle (Gold v2) — v02 (08.09.)
% In Overleaf ablegen als: tables/Chap7-v02/chap7.3table1.tex
% Aufruf im Fließtext: % B-59, 15.09.2026: datierte Ledger-Nachauswertung, Stützung und Reichweite.
% Beleg: thesis-evidence/w2/ledger-repetitions-20260915/ (Originale unverändert).
% M19 (12.09.2026): Leserführung und präzise Ergebnisdarstellung;
% Messwerte und historische Nachweisgrenzen bleiben erhalten.
% ============================================================
% Tabelle 7.3-1: Kernzahlen beider Fälle (Gold v2) — v02 (08.09.)
% In Overleaf ablegen als: tables/Chap7-v02/chap7.3table1.tex
% Aufruf im Fließtext: % B-59, 15.09.2026: datierte Ledger-Nachauswertung, Stützung und Reichweite.
% Beleg: thesis-evidence/w2/ledger-repetitions-20260915/ (Originale unverändert).
% M19 (12.09.2026): Leserführung und präzise Ergebnisdarstellung;
% Messwerte und historische Nachweisgrenzen bleiben erhalten.
% ============================================================
% Tabelle 7.3-1: Kernzahlen beider Fälle (Gold v2) — v02 (08.09.)
% In Overleaf ablegen als: tables/Chap7-v02/chap7.3table1.tex
% Aufruf im Fließtext: \input{tables/Chap7-v02/chap7.3table1}
% Label: t:eval-kernzahlen-v2
% Stil: chap6.1-Konvention.
% HERKUNFT: konsolidierter Auswertungsstand 08.09. =
% runs/w2-nachreview-konsolidierung.json (Nachreview-Runde: 5
% geänderte Ledger-Urteile 025/062/086/105/107); Basis =
% runs/w2-goldv2-eval-ledger-8189ea.json (Interview) +
% runs/w2-goldv2-eval-ledger-e30646.json (meeting-2, unverändert),
% Kopien thesis-evidence/w2/eval/; Nenner = Gold v2 (109/31, Hash ②).
% ============================================================

\begin{table}[htbp]
  \centering
  \small
  \begin{tabular}{>{\raggedright\arraybackslash}p{0.40\textwidth}rr}
    \hline
    & \textbf{Treue-Fall} & \textbf{Stressfall} \\
    & (meeting-2) & (Interview) \\
    \hline
    Quell-Units / Referenzaussagen (Gold v2) & 36 / 31 & 136 / 109 \\
    Gelieferte kanonische Claims & 20 & 47 \\
    Urteile full / partial / none & 26 / 2 / 3 & 54 / 29 / 26 \\
    \hline
    Full Coverage (strict) & 0{,}839 & 0{,}495 \\
    Touched Coverage & 0{,}903 & 0{,}761 \\
    Inhaltliche Korrektheit der Claims & 1{,}0 & 1{,}0 \\
    Semantische Stützung & 1{,}0 & 0{,}957 \\
    \hline
    Lücken gesamt (partial + none) & 5 & 55 \\
    \quad davon mit zuordenbarem Prüfhinweis & 0 & 37 \\
    \quad davon still & 5 & 18 \\
    Silent-Miss-Rate & 1{,}0 & 0{,}327 \\
    Adressierbarkeit (Prüfpotenzial) & 0{,}839 & 0{,}835 \\
    Prüfhinweise gesamt / lückenrelevant & 3 / 0 & 48 / 18 \\
    Signalbezogene Präzision & 0{,}0 & 0{,}375 \\
    \hline
  \end{tabular}
  \caption[Ergebnisqualität der Ledger-Kette vor dem Human-Gate]{Ergebnisqualität der Ledger-Kette (Pre-Gate) gegen den
    Referenzbestand Gold~v2, je Fall ein Lauf gemäß
    dokumentierter Auswahlregel;
    Stressfall-Abdeckung konsolidiert am 08.09.; Stützung am
    15.09. nachgeprüft (45/47, \hyperref[anh:ledger-wiederholungen]{Anhang~D.11}).
    Bewertungsrevisionen betreffen dieselben Ausgaben, keine
    Systemänderung. „Inhaltliche
    Korrektheit`` folgt dem tatsächlichen Urteilsverfahren:
    quellengetragene Inhalte ohne Gold-Gegenstück zählen nicht als
    Fehler (gesonderte \emph{gold\_escapes}-Zählung nur soweit
    erhoben, Anhang~\ref{anh:evaluation}, \hyperref[anh:referenzpruefung]{Abschnitt~D.3}; keine Gold-Precision im engeren Sinn). Adressierbarkeit
    bezeichnet Prüfpotenzial, keine erreichte Abdeckung; jede Lücke
    und jeder Hinweis zählt je Kennzahl nur einmal.}
  \label{t:eval-kernzahlen-v2}
\end{table}

% Label: t:eval-kernzahlen-v2
% Stil: chap6.1-Konvention.
% HERKUNFT: konsolidierter Auswertungsstand 08.09. =
% runs/w2-nachreview-konsolidierung.json (Nachreview-Runde: 5
% geänderte Ledger-Urteile 025/062/086/105/107); Basis =
% runs/w2-goldv2-eval-ledger-8189ea.json (Interview) +
% runs/w2-goldv2-eval-ledger-e30646.json (meeting-2, unverändert),
% Kopien thesis-evidence/w2/eval/; Nenner = Gold v2 (109/31, Hash ②).
% ============================================================

\begin{table}[htbp]
  \centering
  \small
  \begin{tabular}{>{\raggedright\arraybackslash}p{0.40\textwidth}rr}
    \hline
    & \textbf{Treue-Fall} & \textbf{Stressfall} \\
    & (meeting-2) & (Interview) \\
    \hline
    Quell-Units / Referenzaussagen (Gold v2) & 36 / 31 & 136 / 109 \\
    Gelieferte kanonische Claims & 20 & 47 \\
    Urteile full / partial / none & 26 / 2 / 3 & 54 / 29 / 26 \\
    \hline
    Full Coverage (strict) & 0{,}839 & 0{,}495 \\
    Touched Coverage & 0{,}903 & 0{,}761 \\
    Inhaltliche Korrektheit der Claims & 1{,}0 & 1{,}0 \\
    Semantische Stützung & 1{,}0 & 0{,}957 \\
    \hline
    Lücken gesamt (partial + none) & 5 & 55 \\
    \quad davon mit zuordenbarem Prüfhinweis & 0 & 37 \\
    \quad davon still & 5 & 18 \\
    Silent-Miss-Rate & 1{,}0 & 0{,}327 \\
    Adressierbarkeit (Prüfpotenzial) & 0{,}839 & 0{,}835 \\
    Prüfhinweise gesamt / lückenrelevant & 3 / 0 & 48 / 18 \\
    Signalbezogene Präzision & 0{,}0 & 0{,}375 \\
    \hline
  \end{tabular}
  \caption[Ergebnisqualität der Ledger-Kette vor dem Human-Gate]{Ergebnisqualität der Ledger-Kette (Pre-Gate) gegen den
    Referenzbestand Gold~v2, je Fall ein Lauf gemäß
    dokumentierter Auswahlregel;
    Stressfall-Abdeckung konsolidiert am 08.09.; Stützung am
    15.09. nachgeprüft (45/47, \hyperref[anh:ledger-wiederholungen]{Anhang~D.11}).
    Bewertungsrevisionen betreffen dieselben Ausgaben, keine
    Systemänderung. „Inhaltliche
    Korrektheit`` folgt dem tatsächlichen Urteilsverfahren:
    quellengetragene Inhalte ohne Gold-Gegenstück zählen nicht als
    Fehler (gesonderte \emph{gold\_escapes}-Zählung nur soweit
    erhoben, Anhang~\ref{anh:evaluation}, \hyperref[anh:referenzpruefung]{Abschnitt~D.3}; keine Gold-Precision im engeren Sinn). Adressierbarkeit
    bezeichnet Prüfpotenzial, keine erreichte Abdeckung; jede Lücke
    und jeder Hinweis zählt je Kennzahl nur einmal.}
  \label{t:eval-kernzahlen-v2}
\end{table}

% Label: t:eval-kernzahlen-v2
% Stil: chap6.1-Konvention.
% HERKUNFT: konsolidierter Auswertungsstand 08.09. =
% runs/w2-nachreview-konsolidierung.json (Nachreview-Runde: 5
% geänderte Ledger-Urteile 025/062/086/105/107); Basis =
% runs/w2-goldv2-eval-ledger-8189ea.json (Interview) +
% runs/w2-goldv2-eval-ledger-e30646.json (meeting-2, unverändert),
% Kopien thesis-evidence/w2/eval/; Nenner = Gold v2 (109/31, Hash ②).
% ============================================================

\begin{table}[htbp]
  \centering
  \small
  \begin{tabular}{>{\raggedright\arraybackslash}p{0.40\textwidth}rr}
    \hline
    & \textbf{Treue-Fall} & \textbf{Stressfall} \\
    & (meeting-2) & (Interview) \\
    \hline
    Quell-Units / Referenzaussagen (Gold v2) & 36 / 31 & 136 / 109 \\
    Gelieferte kanonische Claims & 20 & 47 \\
    Urteile full / partial / none & 26 / 2 / 3 & 54 / 29 / 26 \\
    \hline
    Full Coverage (strict) & 0{,}839 & 0{,}495 \\
    Touched Coverage & 0{,}903 & 0{,}761 \\
    Inhaltliche Korrektheit der Claims & 1{,}0 & 1{,}0 \\
    Semantische Stützung & 1{,}0 & 0{,}957 \\
    \hline
    Lücken gesamt (partial + none) & 5 & 55 \\
    \quad davon mit zuordenbarem Prüfhinweis & 0 & 37 \\
    \quad davon still & 5 & 18 \\
    Silent-Miss-Rate & 1{,}0 & 0{,}327 \\
    Adressierbarkeit (Prüfpotenzial) & 0{,}839 & 0{,}835 \\
    Prüfhinweise gesamt / lückenrelevant & 3 / 0 & 48 / 18 \\
    Signalbezogene Präzision & 0{,}0 & 0{,}375 \\
    \hline
  \end{tabular}
  \caption[Ergebnisqualität der Ledger-Kette vor dem Human-Gate]{Ergebnisqualität der Ledger-Kette (Pre-Gate) gegen den
    Referenzbestand Gold~v2, je Fall ein Lauf gemäß
    dokumentierter Auswahlregel;
    Stressfall-Abdeckung konsolidiert am 08.09.; Stützung am
    15.09. nachgeprüft (45/47, \hyperref[anh:ledger-wiederholungen]{Anhang~D.11}).
    Bewertungsrevisionen betreffen dieselben Ausgaben, keine
    Systemänderung. „Inhaltliche
    Korrektheit`` folgt dem tatsächlichen Urteilsverfahren:
    quellengetragene Inhalte ohne Gold-Gegenstück zählen nicht als
    Fehler (gesonderte \emph{gold\_escapes}-Zählung nur soweit
    erhoben, Anhang~\ref{anh:evaluation}, \hyperref[anh:referenzpruefung]{Abschnitt~D.3}; keine Gold-Precision im engeren Sinn). Adressierbarkeit
    bezeichnet Prüfpotenzial, keine erreichte Abdeckung; jede Lücke
    und jeder Hinweis zählt je Kennzahl nur einmal.}
  \label{t:eval-kernzahlen-v2}
\end{table}

% Label: t:eval-kernzahlen-v2
% Stil: chap6.1-Konvention.
% HERKUNFT: konsolidierter Auswertungsstand 08.09. =
% runs/w2-nachreview-konsolidierung.json (Nachreview-Runde: 5
% geänderte Ledger-Urteile 025/062/086/105/107); Basis =
% runs/w2-goldv2-eval-ledger-8189ea.json (Interview) +
% runs/w2-goldv2-eval-ledger-e30646.json (meeting-2, unverändert),
% Kopien thesis-evidence/w2/eval/; Nenner = Gold v2 (109/31, Hash ②).
% ============================================================

\begin{table}[htbp]
  \centering
  \small
  \begin{tabular}{>{\raggedright\arraybackslash}p{0.40\textwidth}rr}
    \hline
    & \textbf{Treue-Fall} & \textbf{Stressfall} \\
    & (meeting-2) & (Interview) \\
    \hline
    Quell-Units / Referenzaussagen (Gold v2) & 36 / 31 & 136 / 109 \\
    Gelieferte kanonische Claims & 20 & 47 \\
    Urteile full / partial / none & 26 / 2 / 3 & 54 / 29 / 26 \\
    \hline
    Full Coverage (strict) & 0{,}839 & 0{,}495 \\
    Touched Coverage & 0{,}903 & 0{,}761 \\
    Inhaltliche Korrektheit der Claims & 1{,}0 & 1{,}0 \\
    Semantische Stützung & 1{,}0 & 0{,}957 \\
    \hline
    Lücken gesamt (partial + none) & 5 & 55 \\
    \quad davon mit zuordenbarem Prüfhinweis & 0 & 37 \\
    \quad davon still & 5 & 18 \\
    Silent-Miss-Rate & 1{,}0 & 0{,}327 \\
    Adressierbarkeit (Prüfpotenzial) & 0{,}839 & 0{,}835 \\
    Prüfhinweise gesamt / lückenrelevant & 3 / 0 & 48 / 18 \\
    Signalbezogene Präzision & 0{,}0 & 0{,}375 \\
    \hline
  \end{tabular}
  \caption[Ergebnisqualität der Ledger-Kette vor dem Human-Gate]{Ergebnisqualität der Ledger-Kette (Pre-Gate) gegen den
    Referenzbestand Gold~v2, je Fall ein Lauf gemäß
    dokumentierter Auswahlregel;
    Stressfall-Abdeckung konsolidiert am 08.09.; Stützung am
    15.09. nachgeprüft (45/47, \hyperref[anh:ledger-wiederholungen]{Anhang~D.11}).
    Bewertungsrevisionen betreffen dieselben Ausgaben, keine
    Systemänderung. „Inhaltliche
    Korrektheit`` folgt dem tatsächlichen Urteilsverfahren:
    quellengetragene Inhalte ohne Gold-Gegenstück zählen nicht als
    Fehler (gesonderte \emph{gold\_escapes}-Zählung nur soweit
    erhoben, Anhang~\ref{anh:evaluation}, \hyperref[anh:referenzpruefung]{Abschnitt~D.3}; keine Gold-Precision im engeren Sinn). Adressierbarkeit
    bezeichnet Prüfpotenzial, keine erreichte Abdeckung; jede Lücke
    und jeder Hinweis zählt je Kennzahl nur einmal.}
  \label{t:eval-kernzahlen-v2}
\end{table}
